% Pursuing Stacks (cotes 134-2 to 134-8): archive cote 134-6, site page 79 (PDF page 80), typescript page 397.
% typescript. Draft compared: ps5.tex:5476-5549.
% AI-transcribed by Claude Opus 5.5 (claude-opus-5-5, adaptive thinking, effort high), Message Batches API, batch ps3 (msgbatch_01XdFmuVVBAcgc7EFxdtCqsj), 2026-11-12; not yet checked by eye.
\scanpage{134-6/79}
\mspage{397}

This pairing gives rise, in a symmetric way,  to \deleted{a}\inserted{the} functor (41) (which is
an e quivalence)  \uline{and} to a functor
\[ (41')\xout{)} \qquad A\hat{\ } \hand{\xrightarrow{\approx}} \underline{Hom}_{\hand{!}}(B\hat{\ }, (Sets)) \]
which (it turns out) is none  other (up to canonical isomorphism) than
(41) with B replaces by A (and hence A replaced by B ),  and therefore is
equally an e quivalence.  \uline{Thus the pairing }(42)  \uline{between the two topoi}
\uline{A$\hat{\ }$ and B$\hat{\ }$ has the remarkable property that it defines an equivalence of}
\uline{each of these topoi with the category of} (\uline{set-valued}) \uline{cosheaves on the}
\uline{other}.  I do not know of any other example of a pair of topoi related
in such a remarkable way\xout{s},  which we may express by saying that the two
topoi are ``\uline{dual}'' to each other.

We still have to give an explicit expression for the pairing (42),
plus a convenient notation.  I'll write
\[ (43) \qquad \hand{\delta_A}(F,G) \overset{\hand{def}}{=} F \ast G \qquad \text{for  F in A$\hat{\ }$, G in B$\hat{\ }$ } , \]
and I'll use the canonical equivalence (valid for any pair of small
categories A,B  - not necessarily dual to each other -  and any \uline{U}-catego-%
ry  stable under direct limits),  deduced by twofold application of
prop.3 ,  plus the tautological isomorphism similar to (7) p.385  :
\[ (44) \qquad \underline{Hom}_{\hand{!!}}(\hat{A},\hat{B}; \hand{\mathcal{M}}) \hand{\xrightarrow{\approx}} \underline{Hom}(AxB, \hand{\mathcal{M}}) \quad , \quad F \hand{\mapsto} Fo(\hand{\varepsilon}_{\hand{A}} x \hand{\varepsilon}_B) , \]
which shows that (42) is known up to canonical isomorphism  when we know
it's restriction to the full subcategory AxB = AxA$^{o}$ ,  identifying as
usual an object a of A to it's image in A$\hat{\ }$, and similarly for B.  If
b is an object in A,  we'll denote by b$^{o}$  the same object viewed as an
object of A$^{o}$ = B.  With these conventions (including (43)) we get the
nice formula
\[ (45) \qquad a \ast b^{o} = Hom_A(b,a) \quad ( = Hom_B(a^{o},b^{o}) ) \hand{, \text{for } a,b \text{ in } A,} \]
which has the required symmetry property  -  which,  for general
objects  F in A$\hat{\ }$ and G in B$\hat{\ }$,  can be stated as  a bifunctorial isomor\inserted{phism}
\[ (46) \qquad F \ast G \hand{\simeq} G \ast F \qquad , \]
where the operation $\ast$  in the first member refers to the pair (A,B),  and
in the second to the pair  (B,A)).

From (45)  we easily deduce the more general formula for F$\ast$G , when
either F or G is in A resp. B ,  namely

\margin{79}
